% ============================================================
% Problema 2
% ============================================================
\section{Problema 2: Crucigrama}

\begin{tcolorbox}[enunciado]
Resuelva el crucigrama de abajo.

\begin{figure}[H]
    \centering
    \includegraphics[width=.6\linewidth]{crucigrama.png}
    \caption{Captura de pantalla de la primera aproximación numérica}
\end{figure}

\textbf{Horizontal:}
\begin{itemize}[leftmargin=*, label={}]
    \item \textbf{3.} Tipo de cable utilizado en una conexión switch a switch y router a router.
    \item \textbf{5.} Tipo de mensaje que va dirigido a un solo dispositivo.
    \item \textbf{8.} Capa del modelo OSI donde interactúa el sistema Operativo.
    \item \textbf{11.} El mensaje debe ser entendido por el receptor y no por agentes externos.
    \item \textbf{12.} Equipo sin memoria y siempre es propenso a colisionar.
    \item \textbf{14.} Característica en la arquitectura de una red que hace a que la topología debe ser robusta.
\end{itemize}

\vspace{0.3cm}
\textbf{Vertical:}
\begin{itemize}[leftmargin=*, label={}]
    \item \textbf{1.} Es aquella red que puede transmitir voz, datos y vídeo.
    \item \textbf{2.} Existen dos tipos de topología, Lógica y ...
    \item \textbf{4.} Característica de la seguridad para poder mostrar quien se dice ser.
    \item \textbf{6.} Capa del modelo TCP/IP que engloba las capas 5, 6 y 7 del Modelo OSI.
    \item \textbf{7.} Significa que el servicio siempre debe estar accesible.
    \item \textbf{9.} Clasificación de la red normalmente administrada por algún proveedor.
    \item \textbf{10.} Tiene memoria y por cada interfaz hay un dominio de colisión.
    \item \textbf{13.} Red en la cual no tenemos Broadcast y tienen los siguientes protocolos: Fibra Óptica, HDLC, PPP y Frame Relay.
\end{itemize}
\end{tcolorbox}

\subsection*{Solución}
% --- Escribe aquí la solución --- 